\documentclass[11pt,a4paper]{article}
\usepackage[utf8]{inputenc}
\usepackage[T1]{fontenc}
\usepackage[brazil]{babel}
\usepackage[margin=2.5cm]{geometry}
\usepackage{mathptmx}
\usepackage{enumitem}
\usepackage{booktabs}
\usepackage{tabularx}
\usepackage{array}
\usepackage{hyperref}
\hypersetup{colorlinks=false, pdfborder={0 0 0}}
\usepackage{fancyhdr}

\pagestyle{fancy}
\fancyhf{}
\fancyhead[L]{\small TrackFleet --- Registro das Partes Interessadas}
\fancyhead[R]{\small\thepage}
\renewcommand{\headrulewidth}{0.4pt}

\newcolumntype{L}{>{\raggedright\arraybackslash}X}

\begin{document}

\begin{center}
{\Large\bfseries DOCUMENTO DE PARTES INTERESSADAS DO PROJETO}\\[2pt]
{\large\bfseries TrackFleet --- Gestão de Rastreador de Automóveis}\\[4pt]
Disciplina de Gerenciamento de Projetos $\cdot$ Framework PMBOK\textregistered\ 8\textsuperscript{a} Edição $\cdot$ Domínio: Partes Interessadas\\[2pt]
Integrantes: Enzo Allebrand e Kerlon Ribeiro (dois GPs) $\cdot$ 4\textsuperscript{o} semestre $\cdot$ Data: 27/08
\end{center}

\vspace{2mm}
\hrule
\vspace{4mm}

\noindent\textbf{\large Registro das Partes Interessadas}

\vspace{1mm}
\noindent O registro identifica e analisa os indivíduos e organizações que podem influenciar ou ser impactados pelo projeto, definindo seus interesses, nível de poder e a estratégia inicial de classificação.

\section{Mapeamento e Classificação}

\begin{table}[h]
\centering
\small
\begin{tabularx}{\textwidth}{@{} l L L l l @{}}
\toprule
\textbf{Parte Interessada} & \textbf{Papel e Identificação} & \textbf{Interesses e Expectativas} & \textbf{Poder/Interesse} & \textbf{Estratégia} \\
\midrule
\textbf{Diretoria da Transportadora} & Patrocinador (interno). Aprova recursos, mudanças e encerramento. & Plataforma operante até o fim do semestre, reduzindo custos operacionais e multas. & Alto / Alto & \textbf{Gerenciar de perto} \\
\addlinespace
\textbf{Enzo e Kerlon} & GPs e Equipe de Desenvolvimento. & Concluir o projeto acadêmico com sucesso, mantendo escopo exequível e qualidade técnica. & Alto / Alto & \textbf{Gerenciar de perto} \\
\addlinespace
\textbf{Gestores de Frota} & Usuários finais (externos/clientes). Operam o dashboard no dia a dia. & Sistema fácil de usar, dados em tempo real confiáveis e geração de relatórios úteis. & Baixo / Alto & \textbf{Manter informado} \\
\addlinespace
\textbf{Motoristas} & Usuários impactados. Conduzem os veículos rastreados. & Preocupação com privacidade e receio de uso punitivo do rastreamento. & Baixo / Alto & \textbf{Manter informado} \\
\addlinespace
\textbf{Provedores de Hardware e APIs} & Fornecedores (Google Maps, fabricantes de rastreadores, AWS). & Uso correto dos serviços dentro das cotas e do faturamento das APIs contratadas. & Baixo / Baixo & \textbf{Monitorar} \\
\addlinespace
\textbf{Órgão Regulador (ANPD)} & Entidade externa. Fiscaliza conformidade com a LGPD. & Garantia de privacidade e proteção dos dados de localização dos motoristas. & Alto / Baixo & \textbf{Manter satisfeito} \\
\bottomrule
\end{tabularx}
\end{table}

\section{Análise de Prioridade (Modelo de Saliência)}

\begin{itemize}[itemsep=4pt]
    \item \textbf{Prioridade Máxima (Saliência Alta):} A \textbf{Diretoria da Transportadora} detém \textit{poder} (orçamento e aprovação), \textit{legitimidade} (é dona do negócio) e \textit{urgência} (precisa dos resultados financeiros em curto prazo). Os \textbf{GPs} também se enquadram no centro.
    \item \textbf{Prioridade Média (Legitimidade e Urgência):} \textbf{Gestores de Frota} e \textbf{Motoristas} possuem forte legitimidade e urgência nas suas expectativas, mas não detêm poder direto sobre a execução do projeto.
    \item \textbf{Prioridade de Conformidade (Poder e Legitimidade):} A \textbf{ANPD} (via diretrizes LGPD) possui grande poder de veto legal e legitimidade inquestionável, porém baixa urgência no dia a dia do projeto (atuando como restrição).
\end{itemize}

\end{document}
